
et on va
montrer 6.10 d'abord pour @@M0@@ et @@M1@@ (qui vérifient bien les hypothèses de 6.10). Compte tenu de la remarque faite plus haut sur la filtration de @@M2@@, il suffit pour cela de montrer que :
@@M3@@
Or on a @@M4@@~; cette relation s'écrit donc :
@@M5@@
Posons @@M6@@~; alors @@M7@@, et on est ramené à montrer :
@@M8@@
Or ceci résulte de 6.9.

Pour déduire 6.10 de ce cas particulier, on considère le @@M9@@-@@M10@@-homomorphisme @@M11@@ qui envoie @@M12@@ sur @@M13@@ et @@M14@@ sur @@M15@@ pour tout @@M16@@ et tout @@M17@@ (4.7). Cet homomorphisme est compatible aux @@M18@@-filtrations, et envoie @@M19@@ sur @@M20@@, et @@M21@@ sur @@M22@@, d'où le résultat.

\textbf{Proposition 6.11.} On suppose que @@M23@@ est de caractéristique nulle. Soit @@M24@@ un entier. Alors, le foncteur @@M25@@ est une équivalence de catégories de la catégorie des @@M26@@-algèbres graduées @@M27@@ telles que @@M28@@, @@M29@@ pour @@M30@@, dans la catégorie des @@M31@@-@@M32@@-algèbres augmentées @@M33@@ telles que @@M34@@~; un foncteur quasi-inverse est donné par @@M35@@.

Soit @@M36@@ une @@M37@@-algèbre graduée telle que @@M38@@ pour @@M39@@, munie de la filtration associée à sa graduation~; alors @@M40@@, donc @@M41@@ (6.6.5). Réciproquement, si @@M42@@ est une @@M43@@-@@M44@@-algèbre augmentée telle que @@M45@@, alors @@M46@@, @@M47@@ pour @@M48@@.

On va définir des isomorphismes de foncteurs :
@@M49@@
